\section{Index of notation}
\label{app:notation}

The tables list the symbols that are used in more than one place,
with the place where each is introduced. Letters not listed here have
a meaning local to one proof or one subsection, stated there. A few
letters are deliberately reused in different parts for unrelated
objects, as indicated.

\renewcommand{\arraystretch}{1.12}\setlength{\tabcolsep}{4pt}\raggedbottom

\subsection*{General, and Part~\ref{part:geometry}}

\begin{longtable}{@{}p{.23\linewidth}p{.52\linewidth}p{.19\linewidth}@{}}
\toprule
Symbol & Meaning & Introduced\\
\midrule
\endhead
$\T$, $\av f$, $R_\vartheta$ & circle $\R/2\pi\Z$; angular average;
 rotation by $\vartheta$ & \S\ref{subsec:conventions}\\
$e_j$ & $\sqrt2\cos(j\theta)$ & \S\ref{sec:rosen-reference-packet}\\
$w$, $h$, $E$, $g$ & profile (slope), radial height with $h'=w$,
 $E=1+w^2$, metric $E\,dr^2+r^2d\theta^2$ & \eqref{eq:intro-metric}\\
$z$, $u$, $q$ & height $z=h+qu$; small parameter & \S\ref{subsec:normalization}\\
$\bar g$ & remainder metric $g-dz\otimes dz$ & \S\ref{sec:geometric-closure}\\
$F_\tau$, $\tau$, $\varrho_\vartheta$ & the bending, its phase,
 intrinsic rotation & \eqref{eq:phase-family}\\
$k$, $p$, $D$, $Q$ & $z_r$, $z_\theta/r$, discriminant, $k+Ep_\theta$
 & \eqref{eq:height-jet-notation}\\
$\alpha$ & connection form of the remainder metric & \eqref{eq:conn-one-form}\\
$\mathcal N_h(z)$ & curvature residual & \eqref{eq:cart-curvature-scalar}\\
$\mathcal F_q(w,u)$, $\mathcal A$ & normalized residual
 & \eqref{eq:cart-removable-equation}\\
$a$, $b$, $c$ & coefficients of the linearization
 & \eqref{eq:linear-coefficients}\\
$\mathscr L_w$, $L_j$, $L'_{2,w}$ & linearization at the radial
 height; its $j$th mode; derivative in $w$
 & \eqref{eq:intro-linearization}, \eqref{eq:rad-eq-modal-operator}\\
$K_0$, $K_g$ & $ww'/r$; curvature of $g$
 & \eqref{eq:cart-determinant-identity}, \eqref{eq:radial-curvature}\\
$B$, $\bar c$, $\bar L$ & first residual correction
 & Lemma~\ref{lem:type-completion}\\
$d$, $v$, $c_\pm$ & drift, speed and slopes of characteristics
 & \eqref{eq:completed-characteristics}\\
$\mathfrak F_w$ & connection period & \eqref{eq:cart-mean-path-flux}\\
$T$, $T_m$ & travel angle of a negative component
 & \S\ref{subsec:layout}\\
\bottomrule
\end{longtable}

\subsection*{Part~\ref{part:reference}}

\begin{longtable}{@{}p{.23\linewidth}p{.52\linewidth}p{.19\linewidth}@{}}
\toprule
Symbol & Meaning & Introduced\\
\midrule
\endhead
$X$, $\mathcal P$, $\rho$, $V_0$, $\lambda$ & model coordinate,
 density, potential data, height of a model & Lemma~\ref{lem:liouville}\\
$H_{j,c}$, $H_{j,M}$ & Schr\"odinger operators of the central and
 Morse models & \S\ref{subsec:rad-central-singlepoint}, \S\ref{subsec:morse}\\
$\upsilon$ & $r_X/r$ & \eqref{eq:rad-eq-sc6}\\
$\varpi$, $d\mu$, $B_j$ & $12w/r^2$; $\varpi\,dr$; $-\varpi^{-1}L_j$
 & \eqref{eq:rad-eq-ha2}\\
$A_c$, $a_c$, $c_c$, $\alpha_c$, $\phi_0$ & constants and seed of the
 central model & \eqref{eq:rad-eq-sc1}, \eqref{eq:rad-eq-sc2}\\
$g_0$, $\ell_0$ & temporary border & \S\ref{subsec:rad-central-singlepoint}\\
$A$, $B$, $d$ & Morse parameters (in \S\S\ref{subsec:morse}--\ref{sec:radial-tails})
 & \eqref{eq:rad-eq-rm5}\\
$d_n$, $A_n$, $B_n$ & parameters of cell $n$; $A_n\asymp e^{-n^2}$
 & \S\ref{sec:radial-tails}\\
$\lambda_n$, $\sigma_n$ & height of cell $n$; height ratios
 & \eqref{eq:rad-eq-js2}\\
$S_n$, $H_n$, $q_n$, $\mathcal E_n$ & $\tfrac23n^3$, $e^{-S_n}$,
 $n^{5/2}\log(e+n)$, $e^{q_n}$ & \eqref{eq:rad-eq-endpoint-weights}\\
$M_n$ & $\exp\{1+n^{5/2}\log^2(e+n)\}$ & \eqref{eq:M-weight}\\
$N$, $C_*$, $c_0$ & first cell index; band constant; angle constant
 & \S\ref{sec:radial-tails}\\
$X_c$, $X_L$, $X_R$, $\mathfrak r_n$ & cut and roots of a cell;
 radius factor & \S\ref{subsec:one-cell}, \S\ref{subsec:joins}\\
$\tau_{n,j}$ & action $j\sqrt{A_n}$ of the caps of cell $n$
 & \eqref{eq:cap-action}\\
$\mathsf p_j$, $\Lambda$ & gauge density $wr^{2j^2-1}$;
 $j^2(j^2-1)$ (in \S\ref{sec:radial-tails}) & \eqref{eq:rad-eq-rm13}\\
$I_n$, $J_n$; $\mu_n,\nu_n,R_n^I,R_n$ & mass and resistance
 intervals; masses and resistances & \eqref{eq:mass-resistance-def-new}\\
$a_n$, $y_n$, $b_n$, $c_n$ & node values and coefficients of the
 recurrence & Lemma~\ref{lem:radial-deep-inverse}\\
$\jn_n$, $\mathcal I_n$ & band edges; frequency bands
 & \eqref{eq:rad-eq-rm21}\\
$\zeta_j$, $\Xi_n$, $\Upsilon_n$, $\delta_n$ & state variable on a
 join; height contrast; matching angle
 & \S\ref{subsec:height-selection}\\
$w_*$, $\phi_*$, $\chi_*$, $g_*$, $\ell_*$ & reference profile, seed,
 metric control, border & Prop.~\ref{prop:rad-fixed-theorem}\\
$\mathbb L_{w_*}$, $\mathbb G_{w_*}$ & bordered reference operator and
 its inverse & \eqref{eq:rad-eq-bordered-reference},
 Lemma~\ref{lem:rad-green-completion}\\
$Y_{\gamma,s}$, $\ea$, $\eb$, $\el$ & source norm; exponents of the
 coefficient allowance & \eqref{eq:rad-eq-source-norm},
 Lemma~\ref{lem:inner-inverse}\\
$\vartheta_m$, $\chi_L$ & partition of unity; endpoint step function
 & \S\ref{sec:radial-what}\\
$G_h$, $R_o$, $G$ & inner, outer and matched inverses
 & \eqref{eq:rad-eq-hn9}, \eqref{eq:outer-robin-new}, \eqref{eq:rad-eq-tm11}\\
$\mS$, $\mK$, $\mA$, $\mB$, $D_0$, $D_b$ & matching matrices
 & \S\S\ref{subsec:rad-localized-hilbert-assembly}--\ref{subsec:terminal-matching}\\
$\mathcal K_n$ & transfer bound $e^{Cn^2+C_o\tau_{n,j}}$
 & \eqref{eq:transition-inputs-new}\\
$\mathcal D^{\rm rad}_\gamma$ & weighted graph domain
 & Def.~\ref{def:radial-graph-domain}\\
\bottomrule
\end{longtable}

\subsection*{Part~\ref{part:nonlinear}}

\begin{longtable}{@{}p{.23\linewidth}p{.52\linewidth}p{.19\linewidth}@{}}
\toprule
Symbol & Meaning & Introduced\\
\midrule
\endhead
$x_\partial$, $H$, $E_0$, $M$ & continuous versions of the endpoint
 scales & \eqref{eq:readable-physical-weights}\\
$\ell$, $\varrho$, $\Dr$ & unit-scale radial variable and derivative
 & \S\ref{subsec:spaces}\\
$\beta$, $\epsilon$ & weight exponent $\tfrac34$; gap between
 exponents & \eqref{eq:beta-choice}, \S\ref{subsec:two-exponents}\\
$\mathcal X_\gamma$, $\mathcal X_\gamma^0$, $\mathcal Y_\beta$,
$\mathcal Z_\gamma$, $\mathcal K_\gamma$, $\Qint$ & spaces of heights,
 sources, circle data; integral functional
 & Def.~\ref{def:readable-physical-grading}\\
$\bx=(u,\nu,\kappa,h_L)$, $\Phi_*$ & state; reference height
 $\phi_*e_2$ & \eqref{eq:rewrite-state}\\
$\chi_m$ & travel-angle controls
 & Lemma~\ref{lem:quantitative-optical-controls}\\
$\mathfrak n_s$, $\mathfrak g_{s,m}$; $b_0$, $\epsilon_0$, $d_*$ &
 seminorms of the state; the neighbourhood; fixed loss
 & \S\ref{subsec:neighbourhood}\\
$f$, $h_R$, $\mathfrak e$ & minus return map, right chart, return
 residual & \eqref{eq:full-height-one-chart}\\
$\mathcal M_T$, $\mathcal Q_h$ & derivative of $\mathfrak e$ in the
 chart & \eqref{eq:full-height-chart-derivative}\\
$\widehat L$, $\mathscr E_{\bx}$ & corrected linearization; its error
 & Prop.~\ref{prop:optical-residual-completion}\\
$\sigma$, $\varphi$; $\xi_L$, $\xi_R$, $\xi_c$ & travel coordinates;
 root coordinates & \eqref{eq:fixed-strip-coordinates}\\
$W_i$, $h_i$, $\alpha_i$, $V_i$ & canonical unknown, height factor,
 density, potential & \eqref{eq:fixed-strip-airy-equation}\\
$y$, $\bd$, $P_{\bx}$, $P_*$ & positive coordinate; $y(1-y)$;
 positive canonical operators & \S\ref{subsec:positive-crossing}\\
$\mathscr C_{\bx}$, $\mathscr S_{\bx}$, $\mathbb L_{\bx}$ &
 transformations to canonical form & \S\ref{subsec:canonical-problem}\\
$\mathbb V_{t,s}$ & spaces of canonical sources
 & \eqref{eq:vol-source-spaces}\\
$\mathbb P_{\bx}$, $\bW=(W,a_L,\mu,\lambda)$ & augmented operator and
 unknown & \eqref{eq:vol-abstract-complete-operator}\\
$k_{\bx}$, $s_{\bx}$, $c_*$ & columns of the augmented operator
 & \S\ref{subsec:canonical-problem}\\
$\mathbb G_*$, $\mathbb D_*$, $L_0$ & reference inverse, its domain,
 its loss & Lemma~\ref{lem:two-weights-closed-domain}\\
$|v|_{\mathcal F}$ & $L^1$ norm of the Fourier transform
 & Lemma~\ref{lem:one-high-factor}\\
$\mathfrak g_{s,m}$, $\mathfrak n_{s,m}$ & localized seminorms of the
 deformation and of the state & \S\ref{subsec:neighbourhood}\\
$\delta_n^{\rm root}$ & width of the root collars
 & \eqref{eq:geometry-root-collar-width}\\
$\mathscr R_{\bx}$ & relative operator & \S\ref{subsec:relative}\\
$G_{\bx}$ & inverse at a nearby state
 & Prop.~\ref{prop:rewrite-volume-inverse}\\
$\mathbf R$, $\mathbf V$, $\mathbf Q$ & residual, approximate inverse,
 quadratic error & \S\ref{sec:solving-the-equation}\\
$\mathfrak X$, $\mathfrak Y$ & spaces of increments and residuals
 & \S\ref{sec:solving-the-equation}\\
$C^\circ$, $R^\circ$, $\mathcal D$, $\mathsf J$, $K_{\bx}$, $U_{\bx}$ &
 blocks of the linear system & \eqref{eq:joint-circle-linear-rows}\\
$(\eta,\mu,\varkappa,\psi)$ & increments of $(u,\nu,\kappa,h_L)$
 & \eqref{eq:corrected-coupled-linear-system}\\
$B_\beta$ & space of control increments
 & \S\ref{sec:solving-the-equation}\\
$\mathcal B_{w,U}$, $\mathcal B^\circ_{\mathbf v}$, $\mathcal C$,
$\mathscr P$, $\mathcal U$ & remainders, mean return map, primitive,
 mean of the height & \S\ref{sec:final-proof}\\
\bottomrule
\end{longtable}

\paragraph{Letters with several meanings.}
The following reuses could not be avoided without departing from
standard notation; the context always determines the meaning.
\begin{itemize}
\item $q$ is the small parameter of the height; $q_n$ is the exponent
 of the coefficient allowance.
\item $\lambda$ is the height of a model in Part~\ref{part:reference}
 and the auxiliary scalar of the augmented operator in
 Section~\ref{sec:stable-inverse}.
\item $\mu$ is a mass in Section~\ref{sec:radial-tails}, the measure
 $d\mu=\varpi\,dr$ in
 Section~\ref{sec:radial-reference-inverse}, and the increment of
 $\nu$ in Part~\ref{part:nonlinear}; $\nu_n$ is a mass in
 Section~\ref{sec:radial-tails} and $\nu$ the metric parameter
 elsewhere.
\item $\kappa_i$ is a boundary coefficient in
 Section~\ref{sec:radial-reference-inverse}; $\kappa_m$ are the
 controls in Part~\ref{part:nonlinear}.
\item $\tau$ is the phase of the bending; $\tau_{n,j}$ is an action.
\item $A,B$ are Morse parameters; the matrices $\mA,\mB$ and the
 coefficient $B=E/(rQ)$ are distinguished by type or by context.
\item $\Lambda$ is $j^2(j^2-1)$ in Section~\ref{sec:radial-tails}, a
 matrix of outward derivatives in
 Section~\ref{sec:radial-reference-inverse}, and $|D_\varphi|$ in
 Section~\ref{sec:stable-inverse}.
\item $G$ with a subscript or in blackboard type ($G_h$, $G_o$,
 $\mathbb G_*$, $G_{\bx}$) is an inverse operator; $G(t,\theta)$ is
 the speed of a flow on the circle
 (Lemma~\ref{lem:circle-flow}, Section~\ref{sec:final-proof}); and
 $G(x,\zeta,\theta)$ is the function of the regularized flow at a
 root (Section~\ref{subsec:root-coordinate}).
\item $D$ is the discriminant; $D_u$, $D_\nu$, $D_{\bx}$ are
 derivatives; $D_0$, $D_b$, $D_h$, $D_\kappa$ are matrices or
 operators of Section~\ref{sec:radial-reference-inverse}.
\end{itemize}
